% ============================================================
\section{Technology: Strain-Gauge Shape Reconstruction}
% ============================================================

\subsection{Physical measurement principle}

FlexTail is a \SI{0.9}{\milli\metre}-thick flexible substrate of \ac{PET},
carrying 18 pairs of printed strain gauges distributed along its
length \cite{walkling2025, minktec2025}.
Each pair is arranged symmetrically about the substrate centreline: one gauge on the
left and one on the right of the midline axis.

The measurement principle exploits differential strain.
When the substrate bends --- as the spine curves during movement --- the material on the
convex side experiences tensile stress (elongation) while the concave side experiences
compressive stress simultaneously.
These mechanical stresses alter the electrical resistance of the printed gauge traces
in proportion to the local deformation.

The microcontroller reads the differential resistance between each left--right pair
at high frequency.
The \emph{difference} between the two elements encodes the local bending magnitude
(sagittal curvature radius at that segment).
Crucially, when the loading is asymmetric --- one side stretched more than the other ---
the ratio of the two signals encodes the local axial torsion of the strip about its
longitudinal axis.
Because almost all real-world spinal movements combine sagittal flexion with some
degree of torsion (e.g., asymmetric lifting, golf backswing), this simultaneous
capture of both degrees of freedom from a single sensor pair is a key architectural
advantage.

This direct mechanical coupling between substrate deformation and electrical output
gives the measurement a fundamental advantage over orientation integration:
curvature is measured, not computed by integrating angular velocity.
There is no accumulation of integration noise over time.

\begin{figure}[ht]
	\centering
	\begin{tikzpicture}[>=Stealth, font=\small]
		% ---- Left panel: straight strip ----
		\begin{scope}
			% Strip body
			\fill[gray!20] (0,0) rectangle (3.2, 0.45);
			\draw[thick] (0,0) rectangle (3.2, 0.45);
			% Centreline
			\draw[dashed, gray] (0, 0.225) -- (3.2, 0.225);
			% Left gauge markers (top half, 4 marks)
			\foreach \x in {0.4, 1.0, 1.6, 2.2, 2.8} {
					\draw[minktecBlue, thick] (\x, 0.225) -- (\x, 0.42);
				}
			% Right gauge markers (bottom half, 4 marks)
			\foreach \x in {0.4, 1.0, 1.6, 2.2, 2.8} {
					\draw[red!70!black, thick] (\x, 0.225) -- (\x, 0.03);
				}
			% Labels
			\node[minktecBlue, above] at (1.6, 0.45) {$R_L$};
			\node[red!70!black, below] at (1.6, 0) {$R_R$};
			\node[below=6pt] at (1.6, 0) {\textit{Straight}\quad $\Delta R = 0$};
		\end{scope}

		% Arrow between panels
		\draw[->, thick] (3.5, 0.225) -- (4.3, 0.225)
		node[midway, above, font=\footnotesize] {bend};

		% ---- Right panel: bent strip ----
		\begin{scope}[xshift=4.6cm]
			% Curved strip: convex top edge, concave bottom edge
			% Outer (convex) arc in blue
			\draw[minktecBlue, very thick]
			(0, 0.5) .. controls (1.0, 1.4) and (2.2, 1.4) .. (3.2, 0.5);
			% Inner (concave) arc in red
			\draw[red!70!black, very thick]
			(0, 0.1) .. controls (1.0, 0.9) and (2.2, 0.9) .. (3.2, 0.1);
			% Fill strip area
			\fill[gray!15, opacity=0.6]
			(0, 0.1) .. controls (1.0, 0.9) and (2.2, 0.9) .. (3.2, 0.1) --
			(3.2, 0.5) .. controls (2.2, 1.4) and (1.0, 1.4) .. (0, 0.5) -- cycle;
			% Labels
			\node[minktecBlue, above right] at (2.6, 1.25) {convex (tension)};
			\node[red!70!black, below right] at (2.6, 0.2) {concave (compression)};
			\node[below=6pt, align=center] at (1.6, 0) {$D = R_L - R_R \propto \kappa$};
		\end{scope}
	\end{tikzpicture}
	\caption{Differential strain-gauge principle. On a straight strip (left) the left ($R_L$) and right ($R_R$) gauges are under equal load; their difference $D = 0$. Under bending (right) the convex edge elongates and the concave edge compresses; $D \propto \kappa$ encodes the local curvature~$\kappa$.}
	\label{fig:sensor_principle}
\end{figure}

\subsection{Sacral IMU and Denavit--Hartenberg reconstruction}

The strain-gauge array captures the \emph{relative} shape of the spine --- how each
segment curves relative to its neighbour.
To express this shape in a gravity-referenced coordinate frame, a six-axis \acs{IMU}
(three-axis accelerometer and three-axis gyroscope, \emph{no magnetometer}) is fixed
to the base housing at sacral level (S1).

The software reads the sacral pitch and roll continuously and uses these as the
gravitational anchor for the kinematic chain.
Starting from this anchored base orientation, a series of \ac{DH}
transformations \cite{denavit1955} propagates the local curvature angles upward through the 18
virtual joints, treating each measurement site as an articulating node of a
robot arm.
Successive matrix multiplication yields the absolute 3D Cartesian coordinates (in mm)
of each segment at each time step.

\begin{figure}[ht]
	\centering
	\begin{tikzpicture}[>=Stealth, font=\small,
			joint/.style={circle, draw=minktecNavy, fill=minktecBlue!20, minimum size=7pt, inner sep=0pt},
			lbl/.style={font=\footnotesize, minktecNavy}]

		% Joints: placed along a gentle S-curve to suggest spine shape
		% Coordinates (x, y): base at bottom
		\coordinate (J1) at (0, 0);      % S1 (sacrum)
		\coordinate (J2) at (0.15, 1.0);
		\coordinate (J3) at (0.25, 2.0);
		\coordinate (J4) at (0.20, 3.0);
		\coordinate (J5) at (0.05, 4.0);
		\coordinate (J6) at (-0.15, 5.0);
		\coordinate (J7) at (-0.25, 6.0); % T6

		% Gravity anchor at base
		\draw[thick, ->, gray] (J1) ++(0, -0.8) -- (J1)
		node[midway, right, font=\footnotesize, gray] {gravity};
		\draw[dashed, gray, very thin] (J1) ++(0, -0.8) ++(0.6,0) -- ++(-1.2, 0);

		% Links (spine segments)
		\draw[minktecNavy, thick]
		(J1) -- (J2) -- (J3) -- (J4) -- (J5) -- (J6) -- (J7);

		% Joint nodes
		\foreach \j/\name in {J1/S1, J2/$j_2$, J3/$j_3$, J4/$j_4$, J5/$j_5$, J6/$j_6$, J7/T6} {
				\node[joint] at (\j) {};
				\node[lbl, right=4pt] at (\j) {\name};
			}

		% Angle annotation at J3
		\draw[minktecBlue, ->, thin]
		($(J2)!0.55!(J3)$) arc[start angle=90, end angle=60, radius=0.35]
		node[right, font=\footnotesize, minktecBlue] {$\theta_i$};

		% DH label
		\node[align=left, font=\footnotesize, gray] at (1.5, 3.0)
		{$\mathbf{T}_i = \mathbf{R}(\theta_i)\mathbf{d}_i$};

	\end{tikzpicture}
	\caption{Denavit--Hartenberg kinematic chain. The sacral \acs{IMU} anchors joint~S1 to the gravity frame. Each subsequent joint~$j_i$ through T6 is connected by a local DH transformation~$\mathbf{T}_i$ parameterised by the differential strain angle~$\theta_i$, yielding absolute 3D coordinates of all 18 segments.}
	\label{fig:dh_chain}
\end{figure}

The detailed implementation of the DH reconstruction and the hysteresis-based postural
state filters is described in \cref{sec:software}.

\subsection{Magnetometer-free design}

The deliberate omission of a magnetometer removes a critical failure mode for
in-the-wild use.
Standard IMU fusion algorithms use the magnetometer to correct yaw drift by referencing
magnetic north.
In industrial, clinical, and domestic environments --- steel reinforcement in
buildings, electric motors in care beds, medical equipment, and ordinary
computer monitors --- local field distortions render the compass dysfunctional,
producing unbounded yaw error over multi-hour recordings \cite{madgwick2011}.

Yaw drift is instead managed by the \emph{flexlib} software through kinematic
heuristics and postural state logic described in \cref{sec:software}.
The resulting yaw stability is adequate for the postural and activity-recognition
tasks described in this document.
The system does not claim sub-degree absolute yaw accuracy in free ambulation;
see \cref{sec:limitations} for a quantitative characterisation.

To further minimise signal drift caused by moisture during physical exertion,
MinkTec collaborates with TracXon in the framework of the BMBF project
\emph{ErgoFee} (KMU-innovativ programme) to evaluate alternative pastes and
substrates for screen-printing the conductive traces and to maximise the
mechanical robustness of the printed gauge structures.

\subsection{Textile integration and soft-tissue artefact mitigation}

The sensor strip is sewn into a compression undergarment (Sensor\-shirt), available in
two lengths: 36\,cm (covering the lumbosacral and lower thoracic spine) and 45\,cm
(extending into the upper thoracic region).
The constant, defined contact pressure of the compression knit keeps the strip aligned
with the spinal column throughout movement, minimising the soft-tissue artefacts that
degrade surface-mounted sensors in dynamic tasks.

For scenarios in which a compression garment cannot be worn --- prolonged bed rest,
post-operative monitoring, or spaceflight environments --- the sensor strip can be
fixed directly to the skin using medical-grade adhesive patches, maintaining full
mechanical coupling without external compression.

The total system weight is 30\,g.
The integrated 300\,mAh lithium-ion cell provides up to 30\,h of continuous recording.
The hardware module detaches from the shirt, allowing standard machine laundering
and enabling use in clinical hygiene cycles.
